\input{preamble.tex}
\begin{document}

\begin{titlepage}
\vspace*{18mm}
{\sffamily\small ПРОЕКТНЫЙ ДОКУМЕНТ · ЧЕРНОВИК 0.1 · 10 НОЯБРЯ 2026}\par
\vspace{8mm}
{\sffamily\bfseries\fontsize{26}{30}\selectfont Тенгварская пишущая\\[2pt] машинка ТПМ-2}\par
\vspace{4mm}
{\large Полностью механическое исполнение:\\ салазки знаков как память}\par
\vspace{14mm}
\begin{center}
\setlength\tgem{13.5mm}
\begin{tikzpicture}[x=\tgu,y=\tgu]
\input{gen/fig_trace_q.tex}
\end{tikzpicture}
\end{center}
\vspace{6mm}
\noindent\begin{minipage}{0.86\textwidth}\small
\textbf{О чём документ.} ТПМ-2 — ручная рычажная машинка без мотора и силового вала: всё движение дают палец наборщика и пружинный двигатель каретки. Она печатает тот же оттиск, что и электромеханическая ТПМ-1, но иначе устроена. Продвигающие литеры печатаются как на обычной машинке, от левого края клетки, а каретка проходит ширину буквы после удара. К оси буквы сдвигается не каретка, а сами мёртвые знаки: все техты и подстрочные знаки сидят на двух салазках, которые предыдущая буква отводит назад на нужную величину. Салазки держатся защёлкой до следующей буквы — их положение и есть память машинки. На титуле вверху — оттиск, рассчитанный по модели ТПМ-2, внизу — эталонный набор шрифтом; он совпадает с оттиском ТПМ-1 знак в знак.
\end{minipage}
\vfill
{\small User · продолжение проекта ТПМ-1.}
\end{titlepage}

\tableofcontents
\newpage

%=====================================================================
\section{Можно ли без электричества}

Можно. В ТПМ-1 электричество делало три вещи: давало силу (силовой вал), хранило код последней литеры (регистр) и складывало два перемещения в одно (сумматор). Все три имеют чисто механические замены, и у каждой есть исторический предшественник.

\begin{center}\small
\begin{tabularx}{\linewidth}{@{}>{\raggedright}p{38mm}>{\raggedright}p{46mm}X@{}}
\toprule
Что было в ТПМ-1 & Чем заменено в ТПМ-2 & Предшественник \\
\midrule
силовой вал, цикл 360° & палец наборщика; порядок действий задан ходом клавиши & комптометр: клавиша сама приводит механизм \\
двоичный код на гребёнке, защёлки кода & аналоговый код: зазоры на кодовой пластине, три мерные планки & комптометр: клавиша задаёт точку останова \\
регистр R и сумматор $\Delta = b + a$ & положение двух салазок знаков, удерживаемое защёлкой & — \\
отложенная вторая фаза подачи & не нужна: каретка проходит $w$ после удара, как обычно & Varityper: подача на 2, 3 или 4 единицы по клавише \\
сдвиг каретки на $\pm d$ для подстрочных & отдельные нижние салазки со своим отводом $c = b - d$ & — \\
ходовой винт с шагом 1/64'' & обычная рейка; мелкий шаг даёт передача от колеса спуска & стандартный спуск рычажной машинки \\
\bottomrule
\end{tabularx}
\end{center}

Главная перемена — в распределении ролей. В ТПМ-1 каретка подстраивалась под знаки: вставала на ось буквы и ждала техт. В ТПМ-2 каретка ведёт себя как на любой пропорциональной машинке, а подстраиваются знаки. Это проще механически: салазки лёгкие и ходят не больше чем на 3,6~мм, а каретку с валиком двигать туда-обратно на каждый знак не нужно.

\paragraph{Что остаётся неизменным.} Литеры, их размеры и оси измерены на шрифте Alcarin Tengwar и пересчитаны в единицы $u = \mathrm{em}/16$ так же, как в ТПМ-1. Кегль тот же: em = 1/4 дюйма, $u$ = 1/64 дюйма $\approx$ 0,4~мм. Порядок ударов тот же: сначала тенгва, потом её знаки, в любом режиме.

%=====================================================================
\section{Предшественники}

\begin{description}[style=nextline,leftmargin=0pt]
  \item[Varityper (Hammond), 1930-е.] Машинка с пропорциональной подачей, выросшая из машинки Хаммонда. Колесо спуска проходит две, три или четыре единицы в зависимости от того, какая клавиша нажата. В патентах фирмы есть и связь между молотком и спуском, которая не даёт каретке тронуться раньше удара. Это прямой образец для узла подачи ТПМ-2, только у нас уровней не три, а двенадцать.
  \item[Комптометр, 1887.] Счётная машина без рукоятки: числа вводятся нажатием клавиш, и механизм работает от пальца. Клавиша одновременно задаёт точку останова качающегося сектора и сама же его поворачивает. Ход клавиши поэтому длинный. В поздних моделях был «контроль полного хода»: машина блокировалась, если клавишу не дожали. ТПМ-2 перенимает все три черты.
  \item[Двухсобачечный спуск.] Колесо спуска связано с кареткой через шестерню; подающая и держащая собачки вместе с упором удерживают колесо против пружинного двигателя каретки. На этом стоит спуск почти всех рычажных машинок.
\end{description}

%=====================================================================
\section{Принцип: знаки идут к оси}

\subsection{Геометрия}

Пусть точка удара (неподвижная в корпусе) находится против точки $P$ на бумаге.

\begin{enumerate}
  \item Продвигающая литера отлита \emph{как обычно}: левый край клетки — на оси рычага. Удар кладёт знак так, что его клетка начинается в $P$, а ось техты оказывается в $P + a$.
  \item После удара, при отпускании клавиши, каретка проходит ширину $w$; против точки удара теперь $P + w$.
  \item Надстрочный знак сидит на салазках A, отведённых назад на $b$; его удар приходится в $P + w - b = P + a$ — точно на ось буквы.
  \item Подстрочный знак сидит на салазках B, отведённых на $c = b - d$; его удар — в $P + w - c = P + a + d$, на нижнюю ось.
\end{enumerate}

\begin{figure}[h]
\centering
\setlength\tgem{34mm}
\begin{tikzpicture}[x=\tgu,y=\tgu]
\input{gen/fig_geometry.tex}
\end{tikzpicture}
\caption{tinco, акут и две точки снизу. Красным — точка удара литеры (1) и положение каретки после подачи (2). Удар знака с салазок A (3) приходится на $b$ левее, с салазок B (4) — на $c$ левее. Все размеры — в $u = \mathrm{em}/16$.}
\label{fig:geom}
\end{figure}

\subsection{Код литеры}

Каждая продвигающая литера сообщает машине три числа:

\begin{center}\small
\begin{tabular}{@{}clll@{}}
\toprule
& Что задаёт & Диапазон & Куда идёт \\
\midrule
$w$ & подачу каретки после удара & 4--16\,u (12 уровней) & переставной упор спуска \\
$b$ & отвод салазок A (надстрочные знаки) & 1--8\,u & сектор памяти A \\
$c = b - d$ & отвод салазок B (подстрочные знаки) & 1--9\,u & сектор памяти B \\
\bottomrule
\end{tabular}
\end{center}

Это те же сведения, что в коде ТПМ-1 $(a, b, d)$, только в другой форме: $a = w - b$, $d = b - c$. Пробелы задают $w = 4$ или $6$ и обнуляют обе памяти, так что знак после пробела ляжет на край клетки. Мёртвые знаки кода не несут вовсе.

\subsection{Проверка на модели}

Модель ТПМ-2 (\code{simulate\_mech} в \code{tools/gen\_mech.py}) прогоняется на обоих примерах из \S\ref{sec:examples} вместе с моделью ТПМ-1. Скрипт требует, чтобы каждая литера и каждый знак легли в одну и ту же точку бумаги в обеих моделях, и останавливается при любом расхождении. Расхождений нет: механическая машинка даёт тот же оттиск.

%=====================================================================
\section{Общая схема}

\begin{figure}[h]
\centering
\begin{tikzpicture}[x=1mm,y=1mm]
% platen
\draw[mech,fill=gray!15,rounded corners=2pt] (5,92) rectangle (155,100);
\node[lbl] at (30,96) {валик (на каретке)};
\draw[arr] (125,103) -- node[lbl,above]{подача $w$} (100,103);
% print point
\coordinate (P) at (80,90);
\fill[axisred] (P) circle (1);
\node[lbl,text=axisred,anchor=south west] at (81,90.3) {точка удара};
% segment arc and typebars
\draw[thinmech] (P) ++(205:58) arc (205:300:58);
\foreach \a in {208,214,...,298} {\draw[ink!55] (P) ++(\a:58) -- ++(\a+180:46);}
% sledges
\fill[sledA!25,draw=sledA] (P) ++(302:54) arc (302:318:54) -- ++(318:8) arc (318:302:62) -- cycle;
\fill[sledB!25,draw=sledB] (P) ++(321:54) arc (321:333:54) -- ++(333:8) arc (333:321:62) -- cycle;
\foreach \a in {304,307,...,316} {\draw[sledA] (P) ++(\a:58) -- ++(\a+180:46);}
\foreach \a in {323,327,331} {\draw[sledB] (P) ++(\a:58) -- ++(\a+180:46);}
\node[lbl,text=sledA,anchor=north] at ($(P)+(310:66)$) {салазки A};
\node[lbl,text=sledB,anchor=west] at ($(P)+(327:64)$) {салазки B};
\draw[{Latex[length=1.6mm]}-,sledA,thick] ($(P)+(310:70)+(-8,-4)$) -- ++(10,0) node[lbl,right,text=sledA]{$-b$};
\draw[{Latex[length=1.6mm]}-,sledB,thick] ($(P)+(327:72)+(-8,-5)$) -- ++(10,0) node[lbl,right,text=sledB]{$-c$};
\node[lbl,anchor=north east] at ($(P)+(250:60)$) {продвигающие рычаги (36)};
% type guide
\draw[mech] (76,84) -- (76,88) (84,84) -- (84,88);
\node[lbl,anchor=north] at (80,83.5) {направляющая};
% memory sectors and measuring bails
\draw[mech,fill=white,rounded corners=1pt] (112,20) rectangle (152,36);
\node[lbl,align=center] at (132,28) {сектора памяти A, B\\и защёлки};
\draw[sledA,thick] (132,36) .. controls (132,46) and (122,46) .. ($(P)+(312:63)$);
\draw[sledB,thick] (140,36) .. controls (142,46) and (130,48) .. ($(P)+(330:63)$);
\draw[mech,fill=white,rounded corners=1pt] (8,20) rectangle (50,36);
\node[lbl,align=center] at (29,28) {спуск: колесо, собачка,\\переставной упор $w$};
\draw[thinmech] (29,36) -- (29,90);
\node[lbl,anchor=west,align=left] at (30,60) {к рейке\\каретки};
\draw[mech,fill=white,rounded corners=1pt] (55,4) rectangle (152,14);
\node[lbl] at (103.5,9) {мерные планки W, B, C — под всеми продвигающими клавишами};
\draw[arr] (60,14) -- (40,20);
\draw[arr] (132,14) -- (132,20);
\end{tikzpicture}
\caption{ТПМ-2, вид сверху (условно). 44 литерных рычага сходятся к точке удара; восемь рычагов мёртвых знаков стоят на двух салазках у правого края сегмента. Салазки ходят параллельно валику, и точка удара их рычагов смещается вместе с ними. Сектора памяти ставят салазки, мерные планки задают подачу и память.}
\label{fig:top}
\end{figure}

\paragraph{Из чего состоит.} Остов обычной рычажной машинки: сегмент с 44 рычагами, каретка с валиком и пружинным двигателем, рейка, колесо спуска, клавиша регистра. К нему добавлены:
\begin{itemize}
  \item \textbf{кодовые пластины} на 36 продвигающих клавишах — по два яруса, по три выступа на ярус;
  \item \textbf{три мерные планки} W, B, C под всеми продвигающими клавишами;
  \item \textbf{переставной упор} спуска, который ставит планка W;
  \item \textbf{два сектора памяти} с защёлками, которые ставят планки B и C;
  \item \textbf{салазки A и B} — подвижные части сегмента с рычагами мёртвых знаков;
  \item \textbf{блокировка хода} — не даёт начать вторую продвигающую клавишу, пока не закончилась первая, и не даёт бросить клавишу, не дожав её.
\end{itemize}

%=====================================================================
\section{Узлы}

\subsection{Кодовая пластина и мерные планки}

Под каждой продвигающей клавишей висит кодовая пластина с тремя выступами — по одному над каждой мерной планкой. Планки — длинные качающиеся рамки поперёк всей клавиатуры; пружина держит их в верхнем положении. Код задан \emph{зазором} между выступом и планкой: при полном ходе $T$ выступ с зазором $g$ проходит пустой ход $g$, а затем ведёт планку на $T - g$. Значит, ход планки
\[
  s = T - g = k \cdot v, \qquad g = T - k\,v,
\]
где $v$ — значение ($w$, $b$ или $c$), $k$ — ступень. Это та же идея, что у комптометра: клавиша сама задаёт, насколько повернётся механизм, и сама его поворачивает.

\begin{figure}[h]
\centering
\begin{tikzpicture}[x=1mm,y=1mm]
% key lever
\draw[mech,fill=white] (2,50) rectangle (14,54); \node[lbl,anchor=south] at (8,54.5) {клавиша tinco};
\draw[mech] (14,51) -- (118,51); \fill (118,51) circle (1.2); \node[lbl,anchor=south] at (118,52.5) {ось};
% code plate
\draw[mech,fill=dead!70] (40,51) rectangle (88,45);
\node[lbl] at (64,48) {кодовая пластина (нижний ярус)};
% three pushers with gaps: w=10, b=5, c=4  (scale k=0.7 mm/u, T=12 mm)
\def\T{12}\def\k{0.7}
\foreach \x/\v/\nm/\col in {48/10/W/ink,64/5/B/sledA,80/4/C/sledB} {
  \pgfmathsetmacro\g{\T-\k*\v}
  \draw[mech,fill=white] (\x-2,45) rectangle (\x+2,{45-20+\g});
  \fill[\col!35,draw=\col] (\x-5,{45-20-2.4}) rectangle (\x+5,{45-20});
  \node[lbl,text=\col] at (\x,{45-20-6}) {планка \nm};
  \draw[dim] (\x+6,{45-20+\g}) -- node[lbl,right]{$g$} (\x+6,{45-20});
  \node[lbl,anchor=north] at (\x,{45-20-9}) {$v{=}\v$, $g{=}\pgfmathprintnumber[precision=1]{\g}$\,мм};
}
\draw[dim] (100,51) -- node[lbl,right]{$T = 12$ мм — полный ход} (100,39);
\end{tikzpicture}
\caption{Кодовая пластина клавиши tinco ($w = 10$, $b = 5$, $c = 4$) над тремя мерными планками. Чем больше значение, тем меньше зазор и тем дальше выступ ведёт планку. Числа ориентировочные: ступень $k = 0{,}7$~мм на единицу, полный ход 12~мм; ход планки потом увеличивают рычагом.}
\label{fig:plate}
\end{figure}

\paragraph{Два яруса.} На клавише сидят строчная и прописная литеры с разными кодами (например, tinco и удлинённая tinco). Поэтому выступов на пластине два ряда, а клавиша регистра вместе с сегментом сдвигает и мерные планки, чтобы под выступы попал другой ряд. Это то же решение, что в ТПМ-1; для механики оно естественнее, потому что регистр и так двигает тяжёлый сегмент.

\paragraph{Точность.} Значения целые, поэтому все механизмы, которые ставятся по планкам, имеют фиксатор по целым единицам. Он доводит положение до ближайшей единицы и прощает неточность зазора меньше половины ступени.

\subsection{Спуск с переставным упором}

Колесо спуска связано через шестерню с рейкой каретки; пружинный двигатель каретки всё время стремится повернуть колесо. Один зуб колеса соответствует одной единице хода каретки: мелкость $u$ = 0,4~мм получается передаточным числом, а не мелким зубом, так что ни зуб колеса, ни рейка не мельче обычных.

\begin{figure}[h]
\centering
\begin{tikzpicture}[x=1mm,y=1mm]
\def\R{22}\def\n{40}
\draw[mech,fill=gray!6] (0,0) circle (\R);
\foreach \i in {0,...,39} {\draw[ink!70] ({90+\i*9}:\R) -- ({90+\i*9+4.5}:{\R+2.2}) -- ({90+(\i+1)*9}:\R);}
\fill (0,0) circle (1.2);
\node[lbl] at (0,-5) {колесо спуска};
\draw[arr] (-8,10) arc (130:60:12) node[lbl,above right]{пружина каретки};
% check pawl at 90 deg
\draw[mech,fill=white] (-1.5,{\R+2.4}) -- (1.5,{\R+2.4}) -- (0,{\R-0.5}) -- cycle;
\draw[mech] (0,{\R+2.4}) -- (0,{\R+12}); \node[lbl,anchor=south] at (0,{\R+12}) {держащая собачка Д};
% movable stop at w=10 teeth clockwise (-90 deg)
\def\ang{0}
\draw[thinmech,ink!40] (0,0) -- ({\ang}:{\R+8});
\draw[mech,fill=dead] ({\ang}:{\R-0.5}) -- ({\ang+3}:{\R+3}) -- ({\ang-3}:{\R+3}) -- cycle;
\node[lbl,anchor=west] at ({\ang}:{\R+5}) {переставной упор У};
\draw[{Latex[length=1.8mm]}-{Latex[length=1.8mm]},thick,draw=accent] (90:{\R+5}) arc (90:{\ang+4}:{\R+5});
\node[lbl,text=accent] at (45:{\R+10}) {$w$ зубьев};
% pinion and rack
\draw[mech] (0,0) -- (40,-6);
\draw[mech,fill=gray!15] (40,-6) circle (5); \node[lbl,anchor=north] at (40,-12) {шестерня};
\draw[mech] (28,-0.5) -- (70,-0.5);
\foreach \x in {28,30,...,70} {\draw[ink!60] (\x,-0.5) -- (\x+1,1.3);}
\node[lbl,anchor=south] at (58,1.5) {рейка каретки};
% sequence
\node[lbl,anchor=north west,align=left,text width=56mm] at (34,-18) {
\textbf{При отпускании клавиши:}\\
1. Д выходит из зуба;\\
2. колесо идёт до упора У ($w$ зубьев);\\
3. Д падает в зуб;\\
4. У поднимается и возвращается в ноль.};
\end{tikzpicture}
\caption{Спуск ТПМ-2. Планка W при нажатии отводит упор У вперёд на $w$ зубьев; при отпускании держащая собачка освобождает колесо, и каретка проходит ровно $w$ единиц.}
\label{fig:esc}
\end{figure}

Подача происходит при отпускании клавиши, когда рычаг уже отошёл от валика, — как на обычной машинке и как у Varityper, где спуск связан с молотком, чтобы каретка не трогалась раньше удара.

\subsection{Память: сектора с защёлкой}

Планки B и C поворачивают каждая свой сектор памяти. Сектор тягой соединён со своими салазками. У сектора на дуге зубчатый фиксатор с десятью положениями (0--9) и защёлка.

\begin{figure}[h]
\centering
\begin{tikzpicture}[x=1mm,y=1mm]
\fill (0,0) circle (1.2);
\draw[mech,fill=sledA!10] (0,0) -- (60:30) arc (60:120:30) -- cycle;
\foreach \i in {0,...,9} {\pgfmathsetmacro\a{114-\i*6}
  \draw[ink!70] (\a:30) -- (\a:32);
  \node[lbl,font=\sffamily\tiny] at (\a:35) {\i};}
\draw[mech,fill=white] (0,40) -- (2,32.5) -- (-2,32.5) -- cycle;
\draw[mech] (0,40) -- (0,46); \node[lbl,anchor=south] at (0,46) {защёлка};
\draw[mech] (0,0) -- (-30,-8);
\node[lbl,anchor=east] at (-30,-8) {от планки B};
\draw[mech,sledA,thick] (0,0) -- (40,-6) -- (70,-6);
\node[lbl,anchor=west,text=sledA] at (70,-6) {тяга к салазкам A};
\draw[arr] (-12,22) arc (150:120:18) node[lbl,above left]{пружина к нулю};
\node[lbl,anchor=north west,align=left,text width=70mm] at (-35,-14) {
\textbf{На каждой продвигающей клавише:}\\
1. в начале хода сбрасывающая планка поднимает защёлку, пружина возвращает сектор в ноль;\\
2. планка B ведёт сектор на $b$ делений;\\
3. защёлка падает в ближайшую впадину и держит сектор после отпускания клавиши.};
\end{tikzpicture}
\caption{Сектор памяти A (для салазок B — такой же, его ведёт планка C). Память хранится в положении сектора, а значит, и в положении салазок.}
\label{fig:mem}
\end{figure}

Мёртвые клавиши к мерным и сбрасывающей планкам не касаются, поэтому сколько бы знаков ни бить подряд, салазки стоят на месте. Пробелы ставят оба сектора в ноль, возврат каретки — тоже.

\subsection{Салазки знаков}

Рычаги восьми мёртвых клавиш — пять в салазках A (надстрочные), три в салазках B (подстрочные) — стоят не в общем сегменте, а в двух коротких подвижных участках у его края. Салазки ходят по направляющим параллельно валику; при сдвиге на $s$ все их рычаги сдвигаются на $s$, и точка, куда они бьют, смещается на столько же. У каждых салазок своя узкая направляющая литер, которая едет вместе с ними. Ход — до 9\,u $\approx$ 3,6~мм, масса — несколько рычагов, так что сила на установку мала по сравнению с ударом литеры.

\subsection{Служебные клавиши}

\begin{itemize}
  \item \textbf{Пробелы 4u и 6u} — продвигающие клавиши с постоянной пластиной: $w = 4$ или $6$, $b = c = 0$.
  \item \textbf{Забой} — обычный: поворачивает колесо спуска назад на один зуб, то есть на 1\,u. Память не трогает.
  \item \textbf{«Назад на литеру»} (по желанию) — третий сектор памяти для $w$ и храповой возврат каретки на запомненную ширину. Без этого модуля возврат на букву — $w$ ударов забоя.
  \item \textbf{$\pm u$} — качающийся рычажок: пока он отклонён, обе пары салазок сдвинуты на одну единицу в ту или другую сторону. Ручная поправка оси для редких случаев (подстрочные знаки под alda, нестандартные сочетания).
\end{itemize}

\subsection{Блокировки}

Две блокировки — на тех же принципах, что у рычажных машинок и комптометра.
\begin{itemize}
  \item \textbf{Контроль полного хода.} Начатая продвигающая клавиша должна дойти до конца: храповая гребёнка не даёт ей вернуться с полпути. Иначе мерные планки и память встанут на неполное значение.
  \item \textbf{Очерёдность.} Пока продвигающая клавиша не отпущена (а значит, пока не прошла подача и не легла защёлка), остальные клавиши заперты. Мёртвые клавиши тоже ждут: их удар имеет смысл только после того, как салазки встали.
\end{itemize}

%=====================================================================
\section{Ход клавиши}

Порядок действий в механической машинке задаёт не угол вала, а глубина нажатия.

\begin{figure}[h]
\centering
\begin{tikzpicture}[x=1.05mm,y=1mm]
\draw[ink!50,->] (0,0) -- (118,0) node[lbl,right]{};
\foreach \x/\l in {0/0,20/,40/,60/,80/,100/100\%} {\draw[ink!50] (\x,0) -- (\x,-1.5) node[lbl,below]{\l};}
\node[lbl,anchor=north] at (50,-4) {ход вниз};
\draw[ink!50] (100,0) -- (115,0); \node[lbl,anchor=north] at (108,-4) {отпускание};
\def\row#1#2{\node[lbl,anchor=east] at (-2,#1) {#2};}
\node[lbl,anchor=west,font=\sffamily\footnotesize\bfseries] at (0,46) {продвигающая клавиша};
\row{40}{сброс памяти}  \fill[accent!35] (2,38) rectangle (14,42);
\row{33}{мерные планки W, B, C} \fill[sledA!40] (14,31) rectangle (58,35);
\row{26}{бросок рычага, удар} \fill[axisred!50] (58,24) rectangle (100,28); \fill[axisred] (98,24) rectangle (100,28);
\row{19}{защёлки памяти} \fill[sledB!40] (55,17) rectangle (115,21);
\row{12}{подача каретки $w$} \fill[axisblue!40] (102,10) rectangle (114,14);
\node[lbl,anchor=west,font=\sffamily\footnotesize\bfseries] at (0,58) {мёртвая клавиша};
\row{52}{бросок рычага, удар} \fill[axisred!50] (20,50) rectangle (100,54); \fill[axisred] (98,50) rectangle (100,54);
\end{tikzpicture}
\caption{Порядок действий по ходу клавиши. Продвигающая: сначала сброс памяти, затем замер (W ставит упор спуска, B и C ставят салазки), затем бросок рычага; подача — при отпускании. Мёртвая клавиша только бросает рычаг: её салазки уже стоят.}
\label{fig:stroke}
\end{figure}

Ход продвигающей клавиши длиннее обычного: первая половина уходит на замер. У комптометра, где клавиши тоже приводили механизм, ход был длинным по той же причине. Мёртвым клавишам длинный ход не нужен, но для единообразия под пальцами им можно дать тот же.

%=====================================================================
\section{Типы литер}

В ТПМ-1 все литеры были осевыми. В ТПМ-2 осевые только мёртвые знаки, а продвигающие отлиты как обычно — левым краем клетки на оси рычага. Ось техты у них просто нарисована внутри клетки там, где её дал шрифт; на неё попадут знаки с салазок.

\begin{figure}[h]
\centering
\setlength\tgem{26mm}
\begin{tikzpicture}[x=\tgu,y=\tgu]
\input{gen/fig_slugs_adv.tex}
\end{tikzpicture}
\caption{Продвигающие литеры. Красная линия — ось рычага (левый край клетки), серая — ось техты внутри клетки. Под литерой — её код $w$, $b$, $c$.}
\end{figure}

\begin{figure}[h]
\centering
\setlength\tgem{24mm}
\begin{tikzpicture}[x=\tgu,y=\tgu]
\input{gen/fig_slugs_dead.tex}
\end{tikzpicture}
\caption{Мёртвые литеры на салазках A (зелёные) и B (коричневые). Знак отлит по своему якорю, якорь — на оси рычага. Черта удвоения для lambe поднята внутрь изгиба, как в ТПМ-1.}
\end{figure}

Классы и размеры знаков те же, что в ТПМ-1: один комплект гласных на все тенгвы; черта удвоения и носовая черта — узкая и широкая; особая черта для lambe и alda; комбинированные литеры «ára + гласная».

%=====================================================================
\section{Раскладка}

Механика добавляет одно правило: \textbf{клавиша либо целиком мёртвая, либо целиком продвигающая}. Мёртвый рычаг стоит на салазках, продвигающий — в общем сегменте; один рычаг не может быть и там и там. В ТПМ-1 на регистре гласных стояли «ára + гласная», а на регистре носителей — тильда и черта для lambe. Здесь их пришлось развести:

\begin{itemize}
  \item гласные клавиши (салазки A) на регистре несут тильду и носовые черты; два места на регистре «o» и «u» в резерве;
  \item подстрочные клавиши (салазки B): две точки / точка, черта удвоения узкая / широкая, черта для lambe / метка младшего разряда;
  \item «ára + гласная» переехали на регистр нижнего ряда — к anna, yanta, hyarmen и обоим носителям;
  \item halla и удлинённые anga и ungwe заняли освободившиеся места на регистре ряда цифр.
\end{itemize}

Скрипт проверяет оба правила: что ни одна клавиша не смешивает мёртвую и продвигающую литеры и что каждый знак стоит на своих салазках (надстрочный на A, подстрочный на B).

\begin{figure}[h]
\centering
\begin{tikzpicture}[x=1mm,y=1mm]
\input{gen/keyboard.tex}
\end{tikzpicture}
\caption{Раскладка ТПМ-2. Зелёные клавиши — салазки A, коричневые — салазки B, остальные продвигающие. Внизу клавиши — нижний регистр, вверху фиолетовым — верхний.}
\label{fig:kbd}
\end{figure}

%=====================================================================
\section{Набор: два примера}
\label{sec:examples}

В таблицах: «удар» — точка бумаги, куда пришёлся удар (для литеры — левый край клетки, для знака — ось); «каретка» — где оказалась точка удара после клавиши; «A / B» — отвод салазок. Риски под оттиском: красные — удары литер, зелёные и коричневые — удары с салазок.

\subsection{Квенья: \emph{Elen síla lúmenn' omentielvo}}

\begin{figure}[h]
\centering
\setlength\tgem{12.5mm}
\begin{tikzpicture}[x=\tgu,y=\tgu]
\input{gen/fig_trace_q.tex}
\end{tikzpicture}
\caption{Квенийский режим. Оттиск ТПМ-2 совпадает с оттиском ТПМ-1.}
\end{figure}

{\small\setlength\tgem{5mm}
\begin{longtable}{@{}rlrrc@{}}
\toprule
\# & Клавиша & Удар, u & Каретка, u & A / B, u \\
\midrule\endhead
\tabinput{gen/tab_trace_q.tex}
\bottomrule
\end{longtable}}

\subsection{Синдарин: \emph{pedo mellon a minno}}

\begin{figure}[h]
\centering
\setlength\tgem{14mm}
\begin{tikzpicture}[x=\tgu,y=\tgu]
\input{gen/fig_trace_s.tex}
\end{tikzpicture}
\caption{Синдаринский режим. В «mellon» черта для lambe бьёт с салазок B, отведённых на $c = 5$.}
\end{figure}

{\small\setlength\tgem{5mm}
\begin{longtable}{@{}rlrrc@{}}
\toprule
\# & Клавиша & Удар, u & Каретка, u & A / B, u \\
\midrule\endhead
\tabinput{gen/tab_trace_s.tex}
\bottomrule
\end{longtable}}

%=====================================================================
\section{ТПМ-1 и ТПМ-2}

\begin{center}\small
\begin{tabularx}{\linewidth}{@{}>{\raggedright}p{40mm}>{\raggedright}X>{\raggedright\arraybackslash}X@{}}
\toprule
& ТПМ-1 (электромеханика) & ТПМ-2 (механика) \\
\midrule
источник силы & мотор, силовой вал & палец, пружина каретки \\
к оси идёт & каретка & мёртвые знаки (салазки) \\
код литеры & $(a, b, d)$, двоичный, 12 бит на ярус & $(w, b, c)$, зазоры, 3 выступа на ярус \\
память & регистр защёлок & положение двух секторов и салазок \\
подача & до удара, $\Delta = b_{\text{пред}} + a$ & после удара, $w$ \\
подстрочные знаки & сдвиг каретки $\pm d$ на время удара & отдельные салазки, отвод $c = b - d$ \\
каретка & ходовой винт & рейка и колесо спуска \\
ход клавиши & обычный & длинный (замер + бросок) \\
ограничение раскладки & нет & клавиша либо мёртвая, либо продвигающая \\
перекат при быстром наборе & допускается & запрещён блокировкой \\
оттиск & \multicolumn{2}{c}{одинаковый (проверено моделью)} \\
\bottomrule
\end{tabularx}
\end{center}

%=====================================================================
\section{Варианты}

\subsection*{Вариант на челноке Хаммонда}

Varityper печатал не рычагами, а молотком, бьющим по бумаге против сменного изогнутого челнока с литерами. В таком исполнении все знаки сидят на одном лёгком челноке, и салазки не нужны: при ударе мёртвого знака весь узел «челнок + молоток» отводится назад на $b$ или $c$ из памяти и возвращается после удара. Узел лёгкий, поэтому это, возможно, самое изящное механическое решение; цена — сложность самого механизма Хаммонда и его медленность по сравнению с рычажной машинкой.

\subsection*{Моноширинный вариант}

Если все клетки одной ширины, то $b$ и $c$ у всех литер почти одинаковы: салазки можно закрепить, и машинка превращается в обычную, у которой мёртвые клавиши просто отлиты со смещением. Это самый простой путь к рабочему макету — переделка серийной машинки с заменой литер. Цена та же, что в ТПМ-1: сжатые двойные дуги и разрежённые носители.

\subsection*{Без нижних салазок}

Если отказаться от точной нижней оси и ставить подстрочные знаки по верхней оси, салазки B не нужны — останутся одни салазки A и один сектор памяти. Ошибка подстрочного знака будет равна $d$: у основных тенгв 1\,u ($\approx$ 0,4~мм), у hyarmen — 3\,u. Для черновой машинки приемлемо, для чистовой — нет.

%=====================================================================
\section{Риски и открытые вопросы}

\begin{enumerate}[label=Р\arabic*.]
  \item \textbf{Сила и ход клавиши.} Продвигающая клавиша ведёт три планки, сбрасывает память и бросает рычаг. Ход длиннее обычного, нажим тяжелее; это та же цена, что у комптометра. Нужен расчёт сил и, вероятно, пружины, частично уравновешивающие планки.
  \item \textbf{Салазки в сегменте.} Подвижный участок сегмента и подвижная направляющая — нештатные узлы. Нужно убедиться, что при сдвиге на 3,6~мм рычаги салазок не задевают соседних рычагов общего сегмента; поэтому они стоят у края сегмента, а между салазками и общей частью нужен зазор.
  \item \textbf{Двенадцать уровней подачи.} Varityper обходился тремя (2, 3, 4 единицы). Переставной упор на 4--16 зубьев — тот же принцип, но длиннее ход упора и выше требования к его фиксации.
  \item \textbf{Скорость.} Блокировка очерёдности запрещает перекат: ТПМ-2 медленнее ТПМ-1 и медленнее обычной машинки. Для машинки, на которой печатают по-эльфийски, это, вероятно, приемлемо.
  \item \textbf{Регистр с планками.} Сдвиг мерных планок вместе с сегментом при переключении регистра — дополнительная масса под мизинцем. Альтернатива — пары с одинаковым кодом на клавише; для большинства клавиш их подобрать можно, но не для всех.
  \item \textbf{Память на одну литеру.} Как и в ТПМ-1: знаки ставятся только к последней букве. Добить забытую техту к букве двумя раньше можно только забоем и рычажком $\pm u$.
\end{enumerate}

\paragraph{Нужно решить.} Рычажная схема (основной вариант) или челнок Хаммонда; нужны ли нижние салазки; нужен ли модуль «назад на литеру».

%=====================================================================
\appendix
\section{Коды литер ТПМ-2}

$w$ — подача, $a$ — от края клетки до оси техты (только для справки: в механизм не входит), $b$ — отвод салазок A, $d$ — смещение нижней оси, $c = b - d$ — отвод салазок B. Все величины — в $u = \mathrm{em}/16$; источник — тот же замер шрифта Alcarin Tengwar, что и для ТПМ-1.

{\small\setlength\tgem{5.5mm}
\begin{longtable}{@{}cllrrrrr@{}}
\toprule
& Литера & U+ & $w$ & $a$ & $b$ & $d$ & $c$ \\
\midrule\endhead
\tabinput{gen/tab_codes.tex}
\bottomrule
\end{longtable}}

\section*{Источники}
\begin{itemize}\small
  \item Varityper — \url{https://wikipedia.com/wiki/Varityper}.
  \item US~2581200. Typewriter carriage spacing control (Varityper; спуск на 2, 3 или 4 единицы, связь молотка со спуском) — \url{https://patents.google.com/patent/US2581200}.
  \item US~1857771. Variable feed mechanism for typewriter carriages (Varityper) — \url{https://patents.google.com/patent/US1857771A/en}.
  \item US~2904159. Proportional spacing mechanism for typewriters (колесо спуска, подающая и держащая собачки, упор) — \url{https://patents.google.com/patent/US2904159A/en}.
  \item Adding machine (комптометр: ввод нажатием клавиш, привод от пальца) — \url{https://www.wikipedia.org/wiki/Adding_machine}.
  \item HP Museum: Mechanical adders (комптометр: клавиша задаёт останов и поворачивает сектор; длинный ход; блокировка неполного нажатия) — \url{https://hpmuseum.org/adddir.htm}.
  \item Остальные источники (Alcarin Tengwar, раскладка Смита, Tecendil, IBM Executive) — в документе ТПМ-1.
\end{itemize}

\end{document}
